\section{Three methods}\label{sec:methods}

Throughout, $C$ is a smooth projective geometrically integral curve of genus $g$ over a field $k$, and for a
divisor $D$ we write $L(D)=\{f\in k(C)^\times:\operatorname{div}f+D\ge0\}\cup\{0\}$ and $\ell(D)=\dim L(D)$.
A nonconstant $f\in k(C)$ has degree $\deg f=\deg(f)_\infty$, the degree of its pole divisor, and
$\gon_kC\le d$ if and only if there is an effective divisor $D$ of degree at most $d$ with $\ell(D)\ge2$.

\subsection{Method I: explicit pencils}\label{sec:method-I}
To prove $\gon_\Q C\le d$ we exhibit a rational divisor $D$ of degree $d$ and compute $\ell(D)\ge2$ with
Magma's Riemann--Roch routine. The divisor need not be effective: for $0\ne f\in L(D)$ the divisor
$D+\operatorname{div}f$ is effective, rational and linearly equivalent to $D$. The divisors used are of three kinds:
\begin{enumerate}
\item divisors supported on rational cusps, or on cusps and other low-degree points, whose
 Riemann--Roch spaces we compute from $q$-expansions at all cusps, using Zywina's package
 \cite{ZywinaModular};
\item pullbacks along a forgetful map $X_\Delta(N)\to X_{\Delta'}(M)$ of a pencil of minimal degree on the
 target, whose gonality is known (\cref{lem:pullback});
\item low genus: if $g\ge2$ and $P\in C(\Q)$, then $\ell(K-(g-2)P)=\ell((g-2)P)+1\ge2$ by Riemann--Roch, so
 $\gon_\Q C\le g$; if $g\le1$ and $P\in C(\Q)$, then $\ell(2P)\ge2$ (and $\ell(P)=2$ if $g=0$).
\end{enumerate}
For $X_1(N)$ the upper bounds in Theorem~D are the degrees of the modular units of Derickx and van
Hoeij \cite[Table~1]{DerickxVanHoeij2014}.

\subsection{Method II: finite-field exclusion}\label{sec:method-II}
To prove $\gon_\Q C>d$, we choose a prime $p$ of good reduction (\cref{lem:specialization}) and show that
$C_{\Fp}$ has no nonconstant function of degree at most $d$. For this we use the algorithm of Derickx and Terao
\cite{DerickxTerao}, implemented in their Magma package CurveArith; it works over any finite field $\Fq$ (here
$q=p$) and has two phases.

\emph{Phase 1 (places and expansions).} Let $N_1=\#C(\Fq)$ and $n=\lceil N_1/(q+1)\rceil$. CurveArith
enumerates the places of $C_{\Fq}$ of degree at most $d-\min(n,d-1)$, fixes a basis $\omega_1,\dots,\omega_g$
of $H^0(C,\Omega^1)$, and at each place $P$ computes the first $m_P$ coefficients of the expansions of the
$\omega_i$ in a uniformizer at $P$, as vectors over $\Fq$ ($\deg P$ coordinates per coefficient). The precision is
$m_P=d$ at rational places and $m_P=\lfloor(d-\min(n,d-1))/e\rfloor$ at places of degree $e\ge2$; this suffices
for all divisors of degree $d$ containing at least $\min(n,d-1)$ rational places. For an
effective divisor $D=\sum a_PP$ with $a_P\le m_P$ this gives a $g\times\deg D$ matrix $M_D$, whose columns
are the coefficients of index $<a_P$ at each $P$ in its support.

\begin{lemma}[{\cite{DerickxTerao}}]\label{lem:rank}
$\ell(D)=\deg D+1-\rank M_D$.
\end{lemma}
\begin{proof}
By Riemann--Roch, $\ell(D)=\deg D+1-g+\ell(K-D)$, and $\ell(K-D)=g-\rank M_D$ is the dimension of the space of
holomorphic differentials vanishing on $D$.
\end{proof}

\begin{lemma}[{\cite{DerickxTerao,DerickxVanHoeij2014}, \cite[Lemma~3.1]{NajmanOrlic2024}}]\label{lem:fibres}
If $C_{\Fq}$ has a nonconstant function $f$ of degree at most $d$, it has one whose pole divisor contains at
least $n$ distinct rational places.
\end{lemma}
\begin{proof}
The $N_1$ rational points are distributed among the $q+1$ fibres of $f$ over $\PP^1(\Fq)$, so one fibre
contains at least $n$ of them. Compose $f$ with an automorphism of $\PP^1$ over $\Fq$ moving that fibre to
$\infty$.
\end{proof}

\emph{Phase 2 (search).} CurveArith tests all effective divisors $D$ of degree $d$ containing at least $n$
rational places for $\ell(D)\ge2$ (\cref{lem:rank}). This is enough for every strict intermediate curve in
this paper. For the curves $X_1(N)$ with $41\le N\le50$ there are far too many such divisors, and we use
instead a smaller family of larger divisors, as follows.

\subsection{A smaller search family}\label{sec:phase2}
Let $S=\{P_1,\dots,P_{N_1}\}$ be the rational places, and also write $S=\sum P_i$. Fix an integer $b\ge1$.
For an effective divisor $Q$ with no rational place in its support and an effective divisor $E$ supported on
$S$, put $s=\#\supp E$ and
\begin{equation}\label{eq:budget}
 c(E,Q)=b\,s+\max(n-s,0)+\deg E+\deg Q.
\end{equation}
Let $\mathcal F_b(d)$ be the set of divisors $H=bS+E+Q$ with $\deg Q\le d-n$, $c(E,Q)\le d$, and $E$
maximal with this property for its $Q$ (for the order $E\le E'$).

\begin{proposition}\label{prop:covers}
Suppose $n\ge1$. If $C_{\Fq}$ has a nonconstant function of degree at most $d$, then there is
$H\in\mathcal F_b(d)$ such that $L(H)$ contains a nonconstant function of degree at most $d$.
\end{proposition}
\begin{proof}
By \cref{lem:fibres} there is such an $f$ whose pole divisor $D=A+Q$ has at least $n$ distinct rational
places in its support, where $A$ is supported on $S$ and $Q$ on the other places. Put
$E=\sum_{P\in S}\max(v_P(A)-b,0)P$ and $s=\#\supp E$. A place in $\supp E$ contributes $b+v_P(E)$ to $\deg A$,
and $\supp A$ has at least $\max(n-s,0)$ further places, each contributing at least $1$. Hence
$c(E,Q)\le\deg A+\deg Q=\deg f\le d$, and in particular $\deg Q\le d-n$. Since $v_P(A)\le b+v_P(E)$ for every
$P\in S$, we have $D\le bS+E+Q$. Enlarging $E$ to a maximal $E'$ with $c(E',Q)\le d$ keeps
$D\le bS+E'+Q$, so $f\in L(bS+E'+Q)$ with $bS+E'+Q\in\mathcal F_b(d)$.
\end{proof}

The divisors in $\mathcal F_b(d)$ have degree larger than $d$, so $\ell(H)\ge2$ is not enough: we need the
least degree of a nonconstant function in $L(H)$. It can be read off from the same matrix $M_H$.

\begin{lemma}\label{lem:kernel}
Let $H$ be effective. The space $L(H)/k$ is isomorphic to the space of vectors $x$ with $M_Hx=0$. Under this
isomorphism, the order of the pole of $f$ at $P\in\supp H$ is the largest $j\le v_P(H)$ such that the
coordinates of $x$ belonging to the coefficient of index $j-1$ at $P$ are not all zero.
\end{lemma}
\begin{proof}
For each $P\in\supp H$ let $V_P=\mathcal O_P/\mathfrak m_P^{v_P(H)}$ and $V=\bigoplus V_P$. At each $P$,
CurveArith writes each $\omega_i$ as $g_i\,\eta_P$, where $\eta_P$ is one of the basis differentials that does not vanish at
$P$, and expands $g_i$ in a uniformizer $u_P$; the $i$th row of
$M_H$ is the coordinate vector (over $k$, using a $k$-basis of the residue field at $P$) of the jet of $g_i$ in
$V$, in which, for every $j$, the subspace $\mathfrak m_P^j/\mathfrak m_P^{v_P(H)}$
is cut out by the vanishing of the coordinates of index $<j$; indeed, CurveArith computes the coefficients
one at a time, subtracting a lift of each coefficient and dividing by the uniformizer, and this is
$k$-linear. The pairing $(\phi,h)\mapsto\sum_P\operatorname{Tr}_{k(P)/k}\Res_P(\phi\,h\,\eta_P)$, which is perfect because
$\eta_P$ generates the differentials at $P$, identifies the principal parts $\bigoplus_P\mathfrak m_P^{-v_P(H)}/\mathcal O_P$
with the dual of $V$, a principal part having pole order $\le j$ at $P$ exactly when it annihilates
$\mathfrak m_P^{j}/\mathfrak m_P^{v_P(H)}$. By the residue theorem and Serre duality, a principal part is
the principal part of a function in $L(H)$ (unique up to constants) if and only if it annihilates the image
of $H^0(C,\Omega^1)$ in $V$, that is, if and only if its coordinate vector $x$ satisfies $M_Hx=0$.
\end{proof}

\emph{Precision.} The expansions computed in Phase~1 suffice for every $H\in\mathcal F_b(d)$: at a rational place
$P\in\supp E$ we have $v_P(H)=b+v_P(E)\le bs+\deg E\le d$, which is CurveArith's precision there, and at a place
$P$ of degree $e\ge2$ we have $e\,v_P(H)\le\deg Q\le d-n$, so $v_P(H)\le\lfloor(d-\min(n,d-1))/e\rfloor$, which is
CurveArith's precision there. (The program checks this before searching.)

So for $H\in\mathcal F_b(d)$ we compute the kernel of $M_H$ and check whether some nonzero kernel vector, up
to scalars, gives a pole divisor of degree at most $d$. By \cref{prop:covers}, if this fails for every
$H\in\mathcal F_b(d)$, then $\gon_{\Fq}C>d$. The same argument with $b=0$ and all effective divisors of degree
at most $d$ handles curves without rational points; we did not need this case.

\begin{remark}
\Cref{prop:covers} and \cref{lem:kernel} are an elementary variation on the second phase of
\cite{DerickxTerao}; the first phase is used unchanged. The point is that $\mathcal F_b(d)$ is much smaller
than the set of divisors of degree $d$ with $n$ rational places: the rational part is fixed to $bS$ up to the
small excess $E$, and the kernels of $M_H$, of dimension $\ell(H)-1$, are small, because the divisors $H$ have
degree only slightly above $d$. For a fixed $Q$, the matrices $M_{bS+E+Q}$ share the part belonging to $bS+Q$, so we
eliminate it once and add the part belonging to $E$ incrementally. This variation was first implemented as an option
\texttt{Al:="Cover"} of CurveArith's \texttt{CAHasFunctionOfDegreeAtMost}; for the curves $X_1(N)$ we run
the same search in a short C++ program (\S\ref{sec:computations}).
\end{remark}

\subsection{Method III: lifting}\label{sec:method-III}
When $C_{\Fp}$ has more functions or divisor classes of a given kind than can come from $\Q$, we use the
structure of $J(\Q)$, where $J=\Jac C$, to decide which of them lift. The arguments are of three types, each
explained where it is used:
\begin{enumerate}
\item \emph{rank zero}: if $J(\Q)$ is finite, reduction at an odd good prime is injective on $J(\Q)$, and
 a finite list of classes over $\Fp$ can be compared with explicit rational classes
 (\S\ref{sec:quintic-x0}, \S\ref{sec:quintic-intermediate});
\item \emph{a known positive-rank part}: if $J(\Q)=A(\Q)+T$ for an explicit abelian subvariety $A$ and an
 explicit finite group $T$, the reduction of $J(\Q)$ is contained in an explicit subgroup of $J(\Fp)$
 (the curve $43.132.9.a.2$ in \S\ref{sec:gonality-intermediate}, and Route Q in \S\ref{sec:quintic-x0}, \cref{tab:routeQ});
\item \emph{local rigidity}: if over $\Fp$ the relevant locus agrees, to first order, with a known family
 defined over $\Q$, a formal implicit function argument shows that it agrees over $\Z_p$
 (\cref{lem:local}).
\end{enumerate}
